% GENERATED_BY: oc_core_monograph_translation_draft_pipeline_v1.py
% AI_ASSISTED_TRANSLATION_DRAFT: TRUE
% THEOREM_REVIEW_STATUS: PENDING
% THEOREM_REVIEW_MODE: FORMAL_AI_GATE__CODEX_CLI_CHATGPT
% THEOREM_REVIEW_REVIEWER_ID:
% THEOREM_REVIEW_AUTHORITY_CLASS:
% THEOREM_REVIEWED_AT:
% THEOREM_REVIEW_DOSSIER_REF:
% SOURCE_LANGUAGE: EN
% TARGET_LANGUAGE: RU
% SOURCE_EN_REF: content/19_oc_core_1_3_foundational_consistency.tex
\section{Досье базовой согласованности Core 1.3.3}
\label{sec:oc13-foundational-consistency}

Core 1.3.3 добавляет явный слой согласованности к унаследованному монографическому каркасу.
Его цель не в том, чтобы заменить формальные главы редакторским комментарием,
а в том, чтобы сделать три вещи читаемыми в одном месте:

\begin{itemize}
\item почему нынешняя иерархия фиксируется как \(K_0\) through \(K_{12}\), а
      не произвольным отсечением;
\item какие утверждения в томе являются аксиоматическими, выводными,
      структурными или интерпретативными;
\item как связаны между собой мастер-монография и более узкая журнальная
      экстракция core без бесшумного дрейфа в охвате или нотации.
\end{itemize}

\subsection{Почему иерархия фиксируется как \texorpdfstring{$K_0$--$K_{12}$}{K0-K12}}

Core 1.2 уже содержит достаточно структуры, чтобы потребовать иерархию
больше \(K_{10}\). Смысл не в том, что более высокие метки риторически
привлекательны. Смысл в том, что сам исходный корпус уже различает:

\begin{enumerate}
\item континуумы нижнего уровня, чьи оси, пороги и потоки являются внутренними
      для данного класса моделей;
\item теоретические и метатеоретические континуумы, которые организуют сами
      теории и пространства моделей;
\item дальнейший слой, в котором операторы, металогики и пространства моделей
      становятся явными объектами структурной эволюции;
\item и верхний граничный слой, в котором глобальные условия согласованности
      между такими структурами более высокого порядка сами должны быть
      представлены.
\end{enumerate}

Именно поэтому Core 1.3.3 фиксирует иерархию как \(K_0,\dots,K_{12}\), а не
как \(K_0,\dots,K_{10}\). Остановка на \(K_{10}\) оставляет публичный
исходный корпус внутренне несогласованным, потому что более поздние модули уже
требуют \(K_{11}\) и \(K_{12}\) для формулировки межуровневых переходов и
верхнеграничных условий, о которых заявляет фреймворк.

\subsection{Научная роль верхних уровней}

\paragraph{Уровень \texorpdfstring{$K_{11}$}{K11}.}
\(K_{11}\) — это первый уровень, на котором формальные онтологии, системы
операторов, металогики и пространства моделей сами рассматриваются как
структурированные континуумы, подчиненные условиям допустимости,
преобразования и коллапса.
Этот уровень нужен, когда изучают не только теории, но и эволюцию тех
пространств, в которых существуют теории и операторы построения теорий.

\paragraph{Уровень \texorpdfstring{$K_{12}$}{K12}.}
\(K_{12}\) не вводится как мистический слой завершения.
Это минимально специфицированный верхний слой согласованности, необходимый
когда требуется обсуждать глобальную совместимость, ограниченность и режимы
сбоя по семействам объектов \(K_{11}\).
Его роль структурная, а не триумфалистская: он не позволяет иерархии
делать вид, будто над рекурсивно эволюционирующими метатеоретическими
пространствами нет формальной проблемы.

\subsection{Классы утверждений}

Чтобы том оставался читаемым без размытия научного статуса, Core~1.3.3 использует
в этой главе следующую лексику статуса доказательности. Это локальная
таксономия теорем и охвата утверждений; в разделе
the oc13 theorem discussionroadmapк ней добавляется последующая
операциональная нагрузка контроля для replay, comparator, falsifier и
правил эскалации.

\begin{itemize}
\item \textbf{Axiomatic statements} фиксируют первичные допущения, условия
      допустимости или правила идентичности.
\item \textbf{Derived statements} — это утверждения, несущие теорему, с явной
      зависимостью от поддержки и обязательствами доказательства.
\item \textbf{Empirical statements} требуют автономного source-owned packet,
      canonical proof-routed trace, pinned data route, numerical parameter
      law, observable binding, locked held-out replay procedure, comparator
      audit и явного falsifier до promotion.
\item \textbf{Structural statements} описывают организацию фреймворка,
      иерархию и отображение между модулями или уровнями.
\item \textbf{Interpretive statements} объясняют, что означает формальная
      структура, для какого она читателя и какие последующие проекции она
      может тестировать.
\item \textbf{Frontier statements} могут оставаться в research program, но не
      могут входить в положительную внешнюю науку без formal promotion.
\end{itemize}
Эти локальные классы не предполагаются попарно непересекающимися. Когда
предложение может попасть более чем в один класс, его публичный статус
определяется семантической ролью и функцией доказательства, а не поверхностной
формулировкой:
\begin{enumerate}
\item явное противоречие promoted claim сначала помечает refutation-candidate
      или surface-contradiction case для evidence-state summary audit;
\item эмпирическое содержимое проверяется по packet, data route, replay,
      comparator и falsifier support до того, как доверяют какой-либо
      прозовой метке;
\item выводное теоремное содержимое проверяется по its proof route and
      dependency graph;
\item допущенные первичные предпосылки остаются axiomatic premises;
\item structural mapping и hierarchy claims остаются structural even when
      they mention possible later use;
\item explanatory, audience-facing, or didactic sentences remain interpretive
      unless they assert theorem or empirical force;
\item unpromoted research-program proposals remain frontier statements;
\item pure structural or interpretive content keeps the evidence-state summary-recorded
      support-only row or future-work support row class unless it is explicitly
      attached to evidence-bearing bridge material;
\item bridge-only support row is a stable evidence-support class for rows or
      packets that carry bridge support under explicit scope notes; it is not a
      bounded review-route support queue.
\end{enumerate}
Слова вроде ``frontier'', ``candidate'' или ``projection'' therefore
служат диагностическими сигналами, а не самим правилом assignment класса.

Эти прозаические классы представляют собой локальную таксономию статуса
доказательства, а не замену канонической evidence-state summary-лексики научного класса и
closure-verdict. Release verdicts вроде \texttt{bounded review evidence recorded} и
\texttt{repair-required boundary} остаются отдельным слоем verdict. Соответствие таково:
\begin{itemize}
\item Axiomatic statements являются admitted premises или identity rules. Они
      могут находиться внутри маршрута \occode{bounded review route} как named
      assumptions, но сама аксиома от этого не превращается в proved theorem.
\item Derived statements отображаются в \occode{bounded review route} when the
      support route is closed.
\item Empirical statements отображаются в \occode{bounded review route} only after
      the source-owned packet, proof-routed trace, pinned data route,
      numerical parameter law, observable binding, locked held-out replay
      procedure, comparator audit, and falsifier gates pass.
\item Structural and interpretive statements map to support-only row or
      future-work support row as recorded by evidence-state summary unless the statement is attached
      to an evidence-bearing bridge row.
\item Frontier statements map to future-work support row until a later section or
      release promotes them through the same proof-routed or empirical gates
      used for other positive claims, or records an explicit refutation.
\item Fail-closed bridge, frame, or extension lanes retain the class recorded
      by evidence-state summary; closure failure by itself is not refutation.
\item Statements map to negative-control rejection row only when evidence-state summary records an explicit
      refutation rather than a blocked promotion.
\end{itemize}
Эта глава классифицирует statement classes, которые влияют на статус
теоремы, структурный scope, интерпретацию, empirical promotion и
frontier boundaries.
Operational route discipline рассматривается позже в главах theorem-roadmap и
proof-machinery.
Editorial-bridge language — это не седьмой proof class; это
интерпретативная и release-consistency label для утверждений, которые
соединяют монографию, journal core, dossiers и public package, не добавляя
теоремной силы.
Operational statements — это более поздние exact-match, held-out replay,
statistical-threshold, tail-breach, comparator, falsifier и escalation rules,
связанные с empirical promotion. Structural statements проверяются иначе — по
scope, dependency и source-boundary checks. Оба типа контроля применяются
после того, как здесь уже зафиксированы statement class, numerical parameter
law и observable binding, и ни один из них не становится седьмым proof-status
class.

Journal-core extraction может сужать этот корпус, но не имеет права
смешивать эти классы друг с другом.
Интерпретативное предложение может объяснять теорему; оно не может
подменять теорему.
Структурное замечание может мотивировать проекцию; оно не может
тихо служить доказательством этой проекции.

\subsection{Уровни с доказательной силой и уровни не-утверждения}

Мастер-монография поэтому явно разделяет три слоя:

\begin{enumerate}
\item унаследованный formal backbone, который по-прежнему несет структуру
      core theory;
\item слой ремонта Core 1.3.3, который проясняет witness discipline, fate of
      theorem, chronology и public claim boundary;
\item bounded nonclaim layer, который указывает, что нынешний выпуск пока
      не защищает, даже если более широкая research program все еще это
      преследует.
\end{enumerate}

Это разделение важно для hostile review.
Рецензент должен уметь спросить для любого звучащего сильно предложения:
это axiomatic, derived, empirical, structural, interpretive или frontier?
Если оно derived, чем оно поддержано?
Если оно empirical, какой packet, data route, replay, comparator и
falsifier его поддерживают?
Если оно structural or interpretive, что именно оно не утверждает?
Если оно frontier, где именно ему явно запрещен positive release authority?

\subsection{Связь между мастер-монографией и journal core}

Канонический evidence-state summary является научной authority выпуска, а corpus науки,
принадлежащий исходному источнику, — это locus авторства и свидетельства,
стоящий за этой проекцией. Мастер-монография — это flagship reader-facing
reference, потому что она содержит полную теорию, reader ladder, appendices
и explicit public claim boundary в одном месте.
\emph{Journal Core} — это bounded article artifact. The
\emph{journal-core extraction} — это операция, которая выбирает один
защищаемый theorem-bearing route для editorial inspection, не делая вид, что
вся монография уже сжата в этот более короткий артефакт.

По этой причине правило extraction односторонне:

\[
\text{Master Monograph} \longrightarrow \text{Journal Core},
\]

с явным сохранением нотации, nonclaim boundaries и theorem scope.
Процедура extraction может опускать детали.
Journal Core не может расширять claim.
Публичная граница closure остается scope
\texttt{current bounded science surface}, указанным в
разделе~the public claim-boundary chapter. \texttt{frame-only support row}
closures могут быть \texttt{bounded review evidence recorded} as frame support, но они не становятся
standalone empirical promotions. \texttt{bridge-only support row} rows — это
scope-bounded evidence support, а не public promoted closure; они не
перемаркировываются как proof-routed \texttt{bounded review evidence recorded} rows, если только не
создается отдельный source-owned proof-routed packet, который проходит
канонические gates. \texttt{frontier extension row} и \texttt{negative-control rejection row} lanes остаются
вне public \texttt{bounded review evidence recorded} authority.

\subsection{Последствие для читателя}

Эта глава существует потому, что научная строгость и читаемость здесь не
тянут в противоположные стороны.
Чем явнее досье говорит об иерархии, классе утверждений и границе
доказательства, тем проще это сразу для трех читателей:

\begin{itemize}
\item редактора, которому нужен чистый обзор того, что именно защищается;
\item формального читателя, которому нужны точный scope и дисциплина
      зависимостей;
\item неспециалиста-научного читателя, которому нужна связная точка входа в
      повествование перед погружением в полный технический каркас.
\end{itemize}

\paragraph{Bridge to the next chapter.}
С зафиксированными классами утверждений, границей иерархии и relation от
master к derivative следующая глава сужает фокус еще дальше до одной задачи:
кратчайшего lawful proof-navigation route через несущую theorem spine.

\clearpage
